% PPSN 2026 -- Population Dynamics of IVF-Enhanced SPEA2
\documentclass[runningheads]{llncs}

\usepackage[T1]{fontenc}
\usepackage{graphicx}
\usepackage{amsmath,amssymb}
\usepackage{booktabs}
\usepackage{multirow}
\usepackage{xcolor}
\usepackage{subcaption}
\usepackage{hyperref}
\renewcommand\UrlFont{\color{blue}\rmfamily}
\urlstyle{rm}

% Custom commands
\newcommand{\ivfspea}{IVF/SPEA2}
\newcommand{\deltaigd}{\Delta\text{IGD}}

\begin{document}

\title{Population Dynamics Under Local Intensification:\\
  How an IVF-Inspired Operator Reshapes Search in SPEA2}

\titlerunning{Population Dynamics Under Local Intensification}

% TODO: Update author information
\author{Pedro Sanches\inst{1}
\and Co-Author\inst{1}}

\authorrunning{P. Sanches et al.}

\institute{TODO: Institution}

\maketitle

% ============================================================
\begin{abstract}
Memetic multi-objective evolutionary algorithms augment population-based
search with local intensification operators, yet the \emph{dynamic}
impact of these operators on population structure across generations
remains poorly understood. We present a generation-by-generation
analysis of how an IVF-inspired local search operator alters the
convergence--diversity tradeoff within SPEA2. Using paired same-seed
experiments on 18~benchmark instances spanning four problem families
(ZDT, DTLZ, WFG, MaF), we track six population-level metrics---IGD,
hypervolume, spread, spacing, archive turnover, and non-dominated
fraction---across ${\sim}1{,}000$ generations per run. Our analysis
reveals three distinct behavioural regimes: (i)~early-phase diversity
injection where IVF accelerates exploration, (ii)~mid-phase convergence
amplification on problems with connected Pareto fronts, and
(iii)~late-phase interference on disconnected fronts where IVF disrupts
archive stability. We identify the generation at which IVF's benefit
emerges (crossover generation) and show it correlates with problem
structure. These findings provide mechanistic insight into \emph{when}
and \emph{why} local intensification helps or hurts in archive-based
MOEAs, complementing prior work that established \emph{that} it helps.

\keywords{Multi-objective optimization \and Population dynamics \and
  Memetic algorithms \and SPEA2 \and Local search \and
  Convergence--diversity tradeoff}
\end{abstract}

% ============================================================
\section{Introduction}
\label{sec:intro}

Multi-objective evolutionary algorithms (MOEAs) maintain a population of
candidate solutions that collectively approximates the Pareto front. The
balance between convergence toward the front and diversity along it is
the central design tension in MOEAs~\cite{deb2002fast, zitzler2001spea2}.
Memetic algorithms augment evolutionary search with local intensification
operators, often improving final solution quality~\cite{ishibuchi2015review}.
However, most studies evaluate memetic operators by comparing \emph{final}
metric values (IGD, HV) after a fixed evaluation budget, treating the
algorithm as a black box.

This ``final-snapshot'' evaluation obscures critical questions:
\begin{itemize}
  \item \emph{When} during the search does the local operator begin to
    help (or hurt)?
  \item \emph{How} does it alter the population's diversity--convergence
    trajectory?
  \item Does its impact depend on the \emph{phase} of search (early
    exploration vs.\ late refinement)?
\end{itemize}

We address these questions through a generation-by-generation empirical
study of the IVF-inspired local search operator integrated into
SPEA2~\cite{sanches2026ivfspea2}. The IVF operator selects dissimilar
fathers for crossover and accepts offspring collectively (when average
population fitness improves), creating a distinctive intensification
mechanism whose dynamic footprint has not been characterised.

\paragraph{Contributions.}
\begin{enumerate}
  \item A \textbf{paired same-seed experimental design} that isolates the
    IVF operator's per-generation effect by comparing IVF/SPEA2 and SPEA2
    trajectories from identical initial populations.
  \item \textbf{Six population-level metrics} tracked across ${\sim}1{,}000$
    generations on 18~instances, producing ${\sim}1{,}080$ trajectory pairs.
  \item Identification of \textbf{three behavioural regimes}
    (early-diversity, mid-convergence, late-interference) and the
    \emph{crossover generation} at which IVF's benefit materialises.
  \item \textbf{Mechanistic insight} linking IVF's success to Pareto front
    connectivity and its failure to archive destabilisation on
    disconnected fronts.
\end{enumerate}


% ============================================================
\section{Related Work}
\label{sec:related}

\subsection{Population Dynamics in Evolutionary Algorithms}

The study of population dynamics---how diversity, convergence, and
selection pressure evolve across generations---has a long history in
single-objective EAs~\cite{eiben2003introduction}. In multi-objective
optimisation, population dynamics are more complex because diversity
serves both exploration and approximation purposes. Studies have
analysed archive growth~\cite{knowles2003properties}, selection pressure
in indicator-based MOEAs~\cite{li2014selection}, and
diversity-convergence tradeoffs in decomposition-based
algorithms~\cite{trivedi2017survey}.

However, few studies examine how \emph{memetic operators} specifically
alter population dynamics. Most memetic MOEA papers report final metric
comparisons without trajectory analysis~\cite{ishibuchi2015review}.

\subsection{Memetic Multi-Objective Optimisation}

Memetic algorithms combine evolutionary global search with local
improvement heuristics. In multi-objective settings, local search
operators face the challenge of defining improvement in the presence of
multiple, possibly conflicting objectives~\cite{knowles2005memetic}.
Approaches include scalarisation-based local search, Pareto local
search~\cite{paquete2004stochastic}, and operator-specific
intensification mechanisms.

The IVF-inspired operator~\cite{sanches2026ivfspea2} differs from
classical memetic local search in two key aspects: (i)~it explicitly
selects \emph{dissimilar} fathers for crossover, promoting exploration
within intensification, and (ii)~it uses a \emph{collective} acceptance
criterion (average fitness improvement) rather than individual
replacement. These design choices create a distinctive balance between
intensification and diversification whose dynamic consequences we
investigate.

\subsection{Generation-Level Analysis in MOEAs}

Runtime analysis of MOEAs has focused primarily on theoretical bounds
and anytime performance~\cite{zilberstein1996using}. Empirical
generation-level studies include convergence profile
analysis~\cite{bezerra2017comparing} and online performance
indicators~\cite{li2019quality}. Our work extends this line by providing
paired, same-seed trajectory comparisons that isolate a single
operator's contribution.


% ============================================================
\section{Background: IVF/SPEA2}
\label{sec:background}

\subsection{SPEA2}

The Strength Pareto Evolutionary Algorithm 2
(SPEA2)~\cite{zitzler2001spea2} maintains an archive of size $N$ and
assigns each individual a fitness value based on strength (number of
dominated solutions) and $k$-nearest-neighbour density estimation. The
environmental selection phase retains non-dominated individuals and
applies distance-based truncation when the archive overflows.

\subsection{The IVF Operator}

The IVF-inspired operator~\cite{sanches2026ivfspea2} activates when the
population's mean fitness exceeds a threshold ratio $R$. A fraction $C$
of the population is selected as \emph{mothers}. For each mother, a
\emph{dissimilar father} is chosen by selecting the candidate with
maximum objective-space distance from a top-3 tournament by fitness.
Offspring are generated via SBX crossover ($\eta_c = 20$), and a
fraction $M$ of mothers undergo polynomial mutation on $V$ of their
decision variables.

The collective acceptance criterion continues IVF cycles while the
population's \emph{average} fitness improves, up to a maximum of
$\text{Cycles}$ iterations per generation. This mechanism allows
multiple rounds of intensification when productive and early termination
when not.

\paragraph{Parameters.} $C = 0.12$, $R = 0.225$, $M = 0.3$, $V = 0.1$,
$\text{Cycles} = 2$ (tuned via a three-phase pipeline;
see~\cite{sanches2026ivfspea2}).


% ============================================================
\section{Experimental Design}
\label{sec:experiments}

\subsection{Paired Same-Seed Protocol}

To isolate the IVF operator's effect, we run IVF/SPEA2 and standard
SPEA2 with \textbf{identical random seeds}, ensuring both algorithms
start from the same initial population. This paired design enables
per-generation comparison: any divergence in trajectories is attributable
solely to the IVF operator.

For each of the 18~benchmark instances, we execute 30~independent runs
(distinct seeds), producing 30~paired trajectory comparisons per instance
and $18 \times 30 = 540$ trajectory pairs in total.

\subsection{Benchmark Instances}

We select 18~instances from four standard benchmark families, stratified
by the IVF operator's known effect (established in prior
work~\cite{sanches2026ivfspea2}):

\begin{itemize}
  \item \textbf{Strong positive} (5 instances): WFG1, MaF2, WFG5, DTLZ3,
    ZDT6 (all $M=2$)
  \item \textbf{Moderate positive} (5 instances): WFG4, DTLZ4, ZDT3,
    MaF4 ($M=3$), WFG8 ($M=3$)
  \item \textbf{Neutral} (4 instances): DTLZ1 ($M=3$), MaF5, WFG8,
    WFG6 ($M=3$)
  \item \textbf{Negative} (3 instances): WFG2 ($M=2,3$), WFG9
  \item \textbf{Tri-objective reference} (1 instance): DTLZ5 ($M=3$)
\end{itemize}

This stratification ensures balanced coverage of IVF's behavioural
spectrum.

\subsection{Per-Generation Metrics}

At each generation $g$, we record for both algorithms:

\begin{enumerate}
  \item \textbf{IGD}$(g)$: Inverted Generational Distance to the true
    Pareto front (2{,}000 reference points).
  \item \textbf{HV}$(g)$: Hypervolume, normalised by the true Pareto
    front's extent.
  \item \textbf{Spread}$(g)$: Maximum pairwise Euclidean distance in
    objective space (population diameter).
  \item \textbf{Spacing}$(g)$: Standard deviation of nearest-neighbour
    distances (distribution uniformity).
  \item \textbf{Turnover}$(g)$: Fraction of archive members replaced
    since generation $g{-}1$.
  \item \textbf{ND fraction}$(g)$: Proportion of non-dominated
    individuals in the archive.
\end{enumerate}

Additionally, for IVF/SPEA2, we record whether IVF was activated and the
number of IVF cycles executed at each generation.

\subsection{Implementation}

All experiments use PlatEMO~\cite{tian2017platemo} with $N = 100$,
maxFE $= 100{,}000$. The IVF/SPEA2 variant (IVFSPEA2V2TRACE) and a
corresponding SPEA2 baseline (SPEA2TRACE) are instrumented to record
snapshots after each generation's environmental selection. Source code
and data are available at [repository URL].


% ============================================================
\section{Results}
\label{sec:results}

% TODO: Populate with final results after all 18 cases complete.

\subsection{Convergence Trajectories}

% IGD and HV convergence curves (Fig. dynamics_igd_convergence,
% dynamics_hv_convergence)

\subsection{Diversity Dynamics}

% Spread and spacing profiles (Fig. dynamics_spread)
% Archive turnover (Fig. dynamics_turnover)

\subsection{IVF Activation Patterns}

% Heatmap of IVF activation (Fig. dynamics_ivf_activation_heatmap)
% Correlation between activation frequency and benefit

\subsection{Paired $\Delta$IGD Analysis}

% Per-generation delta IGD (Fig. dynamics_delta_igd)
% Crossover generation analysis (Tab. dynamics_crossover_generations)

\subsection{Behavioural Regime Identification}

% Three-phase model: early, mid, late
% Statistical support via changepoint detection or sliding-window tests


% ============================================================
\section{Discussion}
\label{sec:discussion}

% TODO: Interpret results, connect to problem structure

\subsection{When Does IVF Help?}

% Crossover generation correlates with front connectivity
% Early benefit on multimodal problems (DTLZ3, ZDT6)

\subsection{When Does IVF Hurt?}

% Disconnected fronts (WFG2): IVF disrupts archive stability
% Higher turnover without convergence improvement

\subsection{Mechanistic Interpretation}

% Dissimilar father selection -> diversity injection -> helps when
% front is connected (offspring land on the front)
% Collective acceptance -> sustained intensification -> helps when
% landscape is smooth
% On disconnected fronts: offspring fall between components,
% destabilising the archive

\subsection{Practical Implications}

% Guidelines for when to activate/deactivate IVF
% Potential for adaptive triggering based on turnover/ND-fraction


% ============================================================
\section{Conclusion}
\label{sec:conclusion}

We presented a generation-by-generation empirical analysis of how the
IVF-inspired local search operator reshapes population dynamics within
SPEA2. Using a paired same-seed design across 18~benchmark instances, we
identified three behavioural regimes and showed that IVF's benefit
emerges early in search on problems with connected Pareto fronts but
degrades archive stability on disconnected geometries.

These findings move beyond the ``does it work?'' question to provide
mechanistic understanding of \emph{when} and \emph{why} local
intensification helps in archive-based MOEAs. Future work will explore
adaptive IVF triggering based on online population dynamics indicators
(e.g., turnover rate, ND fraction) to activate intensification only when
it is productive.


% ============================================================
\bibliographystyle{splncs04}
\bibliography{references}

\end{document}
